% Options for packages loaded elsewhere
\PassOptionsToPackage{unicode}{hyperref}
\PassOptionsToPackage{hyphens}{url}
\PassOptionsToPackage{dvipsnames,svgnames,x11names}{xcolor}
%
\documentclass[
  a4paper,
  landscape]{report}

\usepackage{amsmath,amssymb}
\usepackage{iftex}
\ifPDFTeX
  \usepackage[T1]{fontenc}
  \usepackage[utf8]{inputenc}
  \usepackage{textcomp} % provide euro and other symbols
\else % if luatex or xetex
  \usepackage{unicode-math}
  \defaultfontfeatures{Scale=MatchLowercase}
  \defaultfontfeatures[\rmfamily]{Ligatures=TeX,Scale=1}
\fi
\usepackage{lmodern}
\ifPDFTeX\else  
    % xetex/luatex font selection
    \setmainfont[]{Gill Sans}
\fi
% Use upquote if available, for straight quotes in verbatim environments
\IfFileExists{upquote.sty}{\usepackage{upquote}}{}
\IfFileExists{microtype.sty}{% use microtype if available
  \usepackage[]{microtype}
  \UseMicrotypeSet[protrusion]{basicmath} % disable protrusion for tt fonts
}{}
\makeatletter
\@ifundefined{KOMAClassName}{% if non-KOMA class
  \IfFileExists{parskip.sty}{%
    \usepackage{parskip}
  }{% else
    \setlength{\parindent}{0pt}
    \setlength{\parskip}{6pt plus 2pt minus 1pt}}
}{% if KOMA class
  \KOMAoptions{parskip=half}}
\makeatother
\usepackage{xcolor}
\usepackage[margin=20mm,heightrounded]{geometry}
\setlength{\emergencystretch}{3em} % prevent overfull lines
\setcounter{secnumdepth}{5}
% Make \paragraph and \subparagraph free-standing
\makeatletter
\ifx\paragraph\undefined\else
  \let\oldparagraph\paragraph
  \renewcommand{\paragraph}{
    \@ifstar
      \xxxParagraphStar
      \xxxParagraphNoStar
  }
  \newcommand{\xxxParagraphStar}[1]{\oldparagraph*{#1}\mbox{}}
  \newcommand{\xxxParagraphNoStar}[1]{\oldparagraph{#1}\mbox{}}
\fi
\ifx\subparagraph\undefined\else
  \let\oldsubparagraph\subparagraph
  \renewcommand{\subparagraph}{
    \@ifstar
      \xxxSubParagraphStar
      \xxxSubParagraphNoStar
  }
  \newcommand{\xxxSubParagraphStar}[1]{\oldsubparagraph*{#1}\mbox{}}
  \newcommand{\xxxSubParagraphNoStar}[1]{\oldsubparagraph{#1}\mbox{}}
\fi
\makeatother


\providecommand{\tightlist}{%
  \setlength{\itemsep}{0pt}\setlength{\parskip}{0pt}}\usepackage{longtable,booktabs,array}
\usepackage{calc} % for calculating minipage widths
% Correct order of tables after \paragraph or \subparagraph
\usepackage{etoolbox}
\makeatletter
\patchcmd\longtable{\par}{\if@noskipsec\mbox{}\fi\par}{}{}
\makeatother
% Allow footnotes in longtable head/foot
\IfFileExists{footnotehyper.sty}{\usepackage{footnotehyper}}{\usepackage{footnote}}
\makesavenoteenv{longtable}
\usepackage{graphicx}
\makeatletter
\newsavebox\pandoc@box
\newcommand*\pandocbounded[1]{% scales image to fit in text height/width
  \sbox\pandoc@box{#1}%
  \Gscale@div\@tempa{\textheight}{\dimexpr\ht\pandoc@box+\dp\pandoc@box\relax}%
  \Gscale@div\@tempb{\linewidth}{\wd\pandoc@box}%
  \ifdim\@tempb\p@<\@tempa\p@\let\@tempa\@tempb\fi% select the smaller of both
  \ifdim\@tempa\p@<\p@\scalebox{\@tempa}{\usebox\pandoc@box}%
  \else\usebox{\pandoc@box}%
  \fi%
}
% Set default figure placement to htbp
\def\fps@figure{htbp}
\makeatother

\usepackage{booktabs}
\usepackage{longtable}
\usepackage{array}
\usepackage{multirow}
\usepackage{wrapfig}
\usepackage{float}
\usepackage{colortbl}
\usepackage{pdflscape}
\usepackage{tabu}
\usepackage{threeparttable}
\usepackage{threeparttablex}
\usepackage[normalem]{ulem}
\usepackage{makecell}
\usepackage{xcolor}
\makeatletter
\@ifpackageloaded{caption}{}{\usepackage{caption}}
\AtBeginDocument{%
\ifdefined\contentsname
  \renewcommand*\contentsname{Table of contents}
\else
  \newcommand\contentsname{Table of contents}
\fi
\ifdefined\listfigurename
  \renewcommand*\listfigurename{List of Figures}
\else
  \newcommand\listfigurename{List of Figures}
\fi
\ifdefined\listtablename
  \renewcommand*\listtablename{List of Tables}
\else
  \newcommand\listtablename{List of Tables}
\fi
\ifdefined\figurename
  \renewcommand*\figurename{Figure}
\else
  \newcommand\figurename{Figure}
\fi
\ifdefined\tablename
  \renewcommand*\tablename{Table}
\else
  \newcommand\tablename{Table}
\fi
}
\@ifpackageloaded{float}{}{\usepackage{float}}
\floatstyle{ruled}
\@ifundefined{c@chapter}{\newfloat{codelisting}{h}{lop}}{\newfloat{codelisting}{h}{lop}[chapter]}
\floatname{codelisting}{Listing}
\newcommand*\listoflistings{\listof{codelisting}{List of Listings}}
\makeatother
\makeatletter
\makeatother
\makeatletter
\@ifpackageloaded{caption}{}{\usepackage{caption}}
\@ifpackageloaded{subcaption}{}{\usepackage{subcaption}}
\makeatother

\ifLuaTeX
\usepackage[bidi=basic]{babel}
\else
\usepackage[bidi=default]{babel}
\fi
\babelprovide[main,import]{english}
\ifPDFTeX
\else
\babelfont{rm}[]{sans-serif}
\fi
% get rid of language-specific shorthands (see #6817):
\let\LanguageShortHands\languageshorthands
\def\languageshorthands#1{}
\ifLuaTeX
  \usepackage[english]{selnolig} % disable illegal ligatures
\fi
\usepackage{bookmark}

\IfFileExists{xurl.sty}{\usepackage{xurl}}{} % add URL line breaks if available
\urlstyle{same} % disable monospaced font for URLs
\hypersetup{
  pdftitle={Report on EXPO\_00035 ALADDIN Pilot 2 (pos) - 1 cal rep},
  pdflang={en},
  colorlinks=true,
  linkcolor={blue},
  filecolor={Maroon},
  citecolor={Blue},
  urlcolor={Blue},
  pdfcreator={LaTeX via pandoc}}


\title{Report on EXPO\_00035 ALADDIN Pilot 2 (pos) - 1 cal rep}
\author{}
\date{Monday, May 26, 2025}

\begin{document}
\maketitle

\renewcommand*\contentsname{Table of contents}
{
\hypersetup{linkcolor=}
\setcounter{tocdepth}{2}
\tableofcontents
}

\chapter{Project Information}\label{project-information}

\begin{itemize}
\tightlist
\item
  User: Axel Mie
\item
  University: SU
\item
  Project: EXPO\_00035 ALADDIN Pilot 2 (pos) - 1 cal rep
\end{itemize}

1 rep, up to 10 ng/mL

\newpage{}

\chapter{General overview}\label{general-overview}

\begin{itemize}
\tightlist
\item
  Number of samples: 62
\item
  Number of targeted chemicals: 2
\end{itemize}

\begin{longtable*}[t]{lrr}
\toprule
Class & Number of Samples & Percent\\
\midrule
Pool & 2 & 3.2\%\\
Sample & 60 & 96.8\%\\
\bottomrule
\end{longtable*}

\newpage{}

\section{PCA}\label{pca}

\includegraphics[width=1\linewidth,height=\textheight,keepaspectratio]{report-external_files/figure-pdf/pca-1.pdf}

\newpage{}

\section{Chemical summary}\label{chemical-summary}

\begin{longtable}[t]{lrrrrrrrrrlr}
\caption{Summary of Targets}\\
\toprule
Chemical & DF\% & N of sampl. & LOD & LLOQ & Median & Average & Min & Max & R2 & Model & N of points\\
\midrule
\cellcolor{gray!10}{Progesterone} & \cellcolor{gray!10}{100.0} & \cellcolor{gray!10}{62} & \cellcolor{gray!10}{0.02} & \cellcolor{gray!10}{0.02} & \cellcolor{gray!10}{15.99} & \cellcolor{gray!10}{48.46} & \cellcolor{gray!10}{0.16} & \cellcolor{gray!10}{236.59} & \cellcolor{gray!10}{0.994} & \cellcolor{gray!10}{linear-1} & \cellcolor{gray!10}{5}\\
Testosterone & 90.3 & 56 & 0.02 & 0.02 & 0.17 & 0.26 & 0.00 & 2.35 & 1.000 & linear-1 & 5\\
\bottomrule
\multicolumn{12}{l}{\rule{0pt}{1em}\textit{Note: } DF: Detection frequency (percentage of samples above LOD)}\\
\end{longtable}

\textbf{Note}

The unit of concentration is ng/mL.

Limit of detection (LOD) and lower-limit of quantification (LLOQ) are
identified as described below.

\begin{enumerate}
\def\labelenumi{\arabic{enumi}.}
\tightlist
\item
  Beginning with the lowest concentration point (e.g.~\texttt{Cal0}),
  identify the lowest concentration point with a detected signal (or
  peak area) greater than zero.
\item
  Identify LOD as the lowest concentration point at which the average
  signal exceeds three times the standard deviation of the non-zero
  concentration point plus the mean of the non-zero point.
\item
  Identify LLOQ as the lowest concentration point satisfying the
  following two conditions:

  \begin{enumerate}
  \def\labelenumii{\arabic{enumii})}
  \tightlist
  \item
    The average signal exceeds ten times the standard deviation of the
    non-zero concentration point plus the mean of the non-zero point.
  \item
    The relative standard deviation (RSD) of the signal is equal to or
    less than 20\%.
  \end{enumerate}
\end{enumerate}

In order to fit a calibration curve to a meaningful range of
concentration, the fitting range is adjusted. The lower-limit of the
range is chosen as the LLOQ. Concentration points below the limit are
excluded from the calibration curve.

As quite a few concentration points in the calibration curve samples of
a chemical can be too far from the range observed in the samples for
quantification (non-\texttt{CalCurve} and non-\texttt{QC} sample
categories), the upper-limit of the calibration range is set. When there
are sufficient number (≥ 5) of calibration points, that the points that
produce signal after normalization exceeding ten times the
\textbf{highest} signal of the quantification samples are excluded.

If the lower-limit is set so high that fewer than three concentration
points remain, that chemical is excluded from the calibration curve,
because the curve cannot be fitted properly.

After the computation of concentration, the following steps are applied
to the data:

\begin{itemize}
\tightlist
\item
  Any negative concentration is adjusted to zero.
\item
  Any concentration below the lower limit of quantification (LLOQ) is
  adjusted to the half of LLOQ.
\item
  Any concentration below the limit of detection (LOD) is adjusted to
  the 1/4 of LLOQ.
\item
  Any chemicals with values below LLOQ in all samples is removed from
  the data.
\end{itemize}

\newpage{}

\chapter{Invidual chemicals}\label{invidual-chemicals}

\newpage{}

\section{Progesterone}\label{progesterone}

\includegraphics[width=1\linewidth,height=\textheight,keepaspectratio]{report-external_files/figure-pdf/individual-chemicals-1.pdf}

\section{Testosterone}\label{testosterone}

\includegraphics[width=1\linewidth,height=\textheight,keepaspectratio]{report-external_files/figure-pdf/individual-chemicals-2.pdf}




\end{document}
